\documentclass[11pt]{article}
\usepackage[T1]{fontenc}
\usepackage[utf8]{inputenc}
\usepackage{lmodern,microtype}
\usepackage[margin=1.08in]{geometry}
\usepackage{amsmath,amssymb,amsthm,mathtools}
\usepackage{xcolor,hyperref}
\hypersetup{colorlinks=true,linkcolor=blue!55!black,citecolor=blue!55!black,urlcolor=blue!60!black}
\setlength{\parindent}{0pt}
\setlength{\parskip}{0.48em}
\setlength{\emergencystretch}{2em}
\newtheorem{theorem}{Theorem}
\newtheorem{lemma}[theorem]{Lemma}
\newtheorem{proposition}[theorem]{Proposition}
\theoremstyle{remark}
\newtheorem{remark}[theorem]{Remark}

\title{Generic Integer-Order Umbilics of Analytic $W$-Congruences}
\author{[Author Name]}
\date{}

\begin{document}
\maketitle

\begin{abstract}
\textbf{Relation to the problem list.}
This manuscript addresses Question~9 (L9), ``Generic integer
$W$-singularity.''  It gives a full proof within the fixed-$m$ stratum of
real-analytic nondegenerate $W$-congruence germs, using the analytic
coefficient topology.  It makes no claim outside that analytic,
nondegenerate, fixed-parameter setting.

Fix an integer $m\ge1$ and consider the real-analytic nondegenerate
$W$-congruence germs whose normalized umbilic parameter equals $m$.  We
prove that the discriminant is generically of real type $A_m$, with both
real signs allowed.  In hodograph coordinates the $W$-equation becomes
$X\Delta=2\Delta$.  The sole possible nonlinear resonance of $X$ is forced
to vanish, so $X$ is analytically linearizable and the entire discriminant
germ has the form $c_2v^2+\kappa u^{m+1}$.  We then construct an analytic
deformation through genuine $W$-congruences which changes $\kappa$
transversely.  Hence $\kappa\ne0$, equivalently type $A_m$, is open and
dense in the fixed-$m$ stratum.
\end{abstract}

\section{Introduction}

Let a real-analytic line congruence be represented in an affine chart by
\[
 f(x,y)=(x,y,0),
 \qquad
 \xi(x,y)=(P(x,y),Q(x,y),1).
\]
Its extended Weingarten equation is
\begin{align}
 W(P,Q)={}&(Q_y-P_x)(Q_{xx}P_{yy}-P_{xx}Q_{yy})\notag\\
 &-2P_y(P_{xx}Q_{xy}-Q_{xx}P_{xy})
 +2Q_x(P_{xy}Q_{yy}-P_{yy}Q_{xy})=0,\label{eq:W}
\end{align}
and the discriminant is
\begin{equation}\label{eq:delta}
 \delta=(P_x-Q_y)^2+4P_yQ_x.
\end{equation}
Craizer and Garcia proved the generic $A_m$ statement for the first few
integer values and conjectured it for every $m\ge1$ \cite{CraizerGarcia}.
We establish the all-order assertion in the natural fixed-parameter
analytic stratum.

At a normalized nondegenerate umbilic at the origin, the common linear part
has been removed and
\begin{equation}\label{eq:normal}
 p_{20}=p_{02}=q_{11}=0,
 \qquad
 p_{11}=a\ne0,
 \qquad
 q_{02}\ne0,
 \qquad
 p_{11}\ne q_{02},
\end{equation}
where $p_{ij}=\partial_x^i\partial_y^jP(0)$ and similarly for $q_{ij}$.
The normalized parameter is
\begin{equation}\label{eq:m}
 m=-\frac{q_{02}}{p_{11}}.
\end{equation}
Let $\mathcal S_m$ denote the space of real-analytic solutions of
\eqref{eq:W} satisfying the umbilic conditions and
\eqref{eq:normal}--\eqref{eq:m}, modulo the normalizing linear changes.
We use the analytic coefficient topology in which a basic neighborhood
restricts finitely many free Taylor coefficients.

We call a real function germ an unsigned real $A_m$ singularity if it is
analytically right-equivalent to
\[
 \varepsilon_2v^2+\varepsilon_1u^{m+1},
 \qquad \varepsilon_1,\varepsilon_2\in\{\pm1\}.
\]

\begin{theorem}\label{thm:main}
For every integer $m\ge1$, the set of germs in $\mathcal S_m$ whose
discriminant has unsigned real type $A_m$ is open and dense.  More precisely,
near every germ there is an analytic scalar $\kappa$ such that
\[
 \delta\sim_{\mathrm{right}}c_2v^2+\kappa u^{m+1},
 \qquad c_2\ne0,
\]
and the exceptional locus $\kappa=0$ is locally a smooth hypersurface.
\end{theorem}

\section{Hodograph reduction}

Assume first that $m>1$.  Set
\begin{equation}\label{eq:AB}
 A=P_y,
 \qquad
 B=\frac{P_x-Q_y}{2}.
\end{equation}
At the normalized umbilic,
\[
 A=ax+O(2),
 \qquad
 B=\frac{m+1}{2}ay+O(2),
\]
so $(A,B)$ are analytic local coordinates.  Define analytic functions
$F,G$ by
\begin{equation}\label{eq:FG}
 -Q_x=F(A,B),
 \qquad
 \frac{P_x+Q_y}{2}=G(A,B).
\end{equation}
Then
\[
 P_x=B+G,
 \quad P_y=A,
 \quad Q_x=-F,
 \quad Q_y=G-B.
\]
Put
\begin{align}
 E&=(1-G_B)F-(1+G_B)AF_A+(AG_A-B)F_B-2BG_A,
 \label{eq:E}\\
 X&=A(1+G_B)\partial_A+(B-AG_A)\partial_B,
 \label{eq:X}\\
 \Delta&=B^2-AF.\label{eq:Delta}
\end{align}

\begin{lemma}\label{lem:identities}
If $J=\det(\partial(A,B)/\partial(x,y))$, then
\begin{equation}\label{eq:identities}
 W=2JE,
 \qquad
 (X-2)\Delta=AE,
 \qquad
 \delta=4\Delta.
\end{equation}
For a $W$-congruence germ,
\begin{equation}\label{eq:eigen}
 X\Delta=2\Delta.
\end{equation}
The linear terms are
\begin{equation}\label{eq:linear}
 X=X_0+O(2),
 \qquad
 X_0=\alpha A\partial_A+B\partial_B,
 \qquad
 \alpha=\frac{2}{m+1},
 \qquad
 \Delta=B^2+O(3).
\end{equation}
\end{lemma}

\begin{proof}
Substituting the four derivative relations following \eqref{eq:FG} into
\eqref{eq:W} and \eqref{eq:delta} gives the first and third identities.
A direct application of \eqref{eq:X} to $B^2-AF$ gives the second.
Since $J(0)=\frac{m+1}{2}a^2\ne0$, $W=0$ implies $E=0$ and hence
\eqref{eq:eigen}.  The normalized two-jet gives \eqref{eq:linear}.
\end{proof}

For $m=1$, the same identities hold and $(A,B)$ are still coordinates,
but $F$ may have a linear $A$-term.  We treat that case separately below.

\section{Analytic normal form}

\begin{proposition}\label{prop:linearization}
For every integer $m\ge2$, the vector field $X$ is analytically linearizable:
there are analytic coordinates $(u,v)$ in which
\begin{equation}\label{eq:linearized}
 X=\alpha u\partial_u+v\partial_v.
\end{equation}
\end{proposition}

\begin{proof}
The eigenvalues $\alpha=2/(m+1)$ and $1$ are positive, so convergent
Poincar\'{e}--Dulac normalization applies \cite{Arnold}.  For a nonlinear
monomial vector field,
\begin{align*}
 [X_0,A^iB^j\partial_A]
 &=\bigl(\alpha(i-1)+j\bigr)A^iB^j\partial_A,\\
 [X_0,A^iB^j\partial_B]
 &=\bigl(\alpha i+j-1\bigr)A^iB^j\partial_B.
\end{align*}
There is no nonlinear $A$-component resonance.  A $B$-component resonance
exists only when $m$ is odd, and then the only possible resonant term is
\[
 cA^r\partial_B,
 \qquad r=\frac{m+1}{2}.
\]
Normalize through degree $r$.  At degree $r+1$, equation
\eqref{eq:eigen} and the quadratic term $\Delta_2=B^2$ give
\begin{equation}\label{eq:kill}
 (X_0-2)\Delta_{r+1}+2cA^rB=0.
\end{equation}
Because $\alpha r+1=2$, the monomial $A^rB$ lies in the kernel of the
scalar homological operator $X_0-2$.  Taking its coefficient in
\eqref{eq:kill} forces $c=0$.  Thus every nonlinear resonance vanishes,
and the convergent normal form is exactly \eqref{eq:linearized}.
\end{proof}

\begin{proposition}\label{prop:classification}
For $m\ge2$, the entire analytic germ $\Delta$ is right-equivalent to
\begin{equation}\label{eq:split}
 c_2v^2+\kappa u^{m+1},
 \qquad c_2\ne0.
\end{equation}
It has type $A_m$ if and only if $\kappa\ne0$; if $\kappa=0$, it has type
$A_\infty$.
\end{proposition}

\begin{proof}
Expand $\Delta$ in the coordinates of \eqref{eq:linearized}.  Equation
\eqref{eq:eigen} allows a monomial $u^iv^j$ only when
\begin{equation}\label{eq:weight}
 \alpha i+j=2.
\end{equation}
The nonnegative integer solutions are $(m+1,0)$ and $(0,2)$, together with
$((m+1)/2,1)$ when $m$ is odd.  Therefore
\[
 \Delta=
 \begin{cases}
  c_2v^2+c_0u^{m+1},&m\text{ even},\\[1mm]
  c_2v^2+c_1u^{(m+1)/2}v+c_0u^{m+1},&m\text{ odd},
 \end{cases}
\]
with $c_2\ne0$.  In the odd case, completing the square by
$v\mapsto v+(c_1/2c_2)u^{(m+1)/2}$ gives \eqref{eq:split}, where
\[
 \kappa=c_0-\frac{c_1^2}{4c_2}.
\]
The even case has $\kappa=c_0$.  Since $\delta=4\Delta$, the singularity
classification follows.
\end{proof}

When $m=1$, the source normalization gives directly
\begin{equation}\label{eq:m1}
 \delta_2=4a\bigl(ay^2+q_{20}x^2\bigr).
\end{equation}
Thus $q_{20}\ne0$ is precisely the unsigned real $A_1$ condition, while
$q_{20}=0$ is the rank-one exceptional locus.

\section{An analytic transverse deformation}

We now show that $\kappa$ is not constrained by the analytic $W$-equation.
Keep $G$ fixed and put
\begin{equation}\label{eq:lambda}
 \lambda=1-G_B,
 \qquad
 \lambda(0)=\frac{2m}{m+1}=m\alpha.
\end{equation}
Use the linearizing coordinates of Proposition~\ref{prop:linearization},
so $X=\alpha u\partial_u+v\partial_v$.  The cohomological equation
\begin{equation}\label{eq:cohomology}
 X\phi=\lambda-\lambda(0)
\end{equation}
has an analytic solution: every nonconstant monomial is divided by the
strictly positive weight $\alpha i+j$.  Hence
\begin{equation}\label{eq:H}
 H=e^\phi u^m
\end{equation}
is analytic and satisfies
\begin{equation}\label{eq:Hhomogeneous}
 XH=(1-G_B)H.
\end{equation}
In the original hodograph coordinates, $H$ has leading term a nonzero
multiple of $A^m$.

Since
\[
 E(F,G)=(1-G_B)F-XF-2BG_A,
\]
the deformation
\begin{equation}\label{eq:deform}
 F_t=F+tH,
 \qquad
 G_t=G
\end{equation}
satisfies $E(F_t,G_t)=0$.  Moreover,
\[
 \Delta_t=\Delta-tAH,
 \qquad
 X(AH)=2AH.
\]
For $m\ge2$, $AH$ has ordinary order $m+1$.  Among the weight-two
monomials in \eqref{eq:weight}, only $u^{m+1}$ has that order.  Therefore
\begin{equation}\label{eq:kappat}
 AH=Cu^{m+1},
 \qquad C\ne0,
 \qquad
 \kappa_t=\kappa-tC.
\end{equation}

It remains to verify that \eqref{eq:deform} comes from genuine analytic
functions $P_t,Q_t$.  Let $x=x(A,B)$ and $y=y(A,B)$, and define
\[
 \omega_P=(B+G)\,dx+A\,dy,
 \qquad
 \omega_Q=-F_t\,dx+(G-B)\,dy.
\]
The closedness conditions are
\begin{align}
 (1+G_B)x_A-G_Ax_B-y_B&=0,\label{eq:CK1}\\
 (F_t)_Bx_A-(F_t)_Ax_B+(1-G_B)y_A+G_Ay_B&=0.
 \label{eq:CK2}
\end{align}
As a system for $(x_A,y_A)$, its determinant at the origin is
\begin{equation}\label{eq:det}
 (1+G_B(0))(1-G_B(0))
 =\frac{4m}{(m+1)^2}>0.
\end{equation}
Choose analytic initial data on $A=0$ from the inverse hodograph map of the
original germ.  Cauchy--Kowalevski applied to
\eqref{eq:CK1}--\eqref{eq:CK2} gives analytic $x_t,y_t$ with nonvanishing
Jacobian for small $t$.  Integrating the closed forms gives analytic
$P_t,Q_t$.  Lemma~\ref{lem:identities} and $E(F_t,G)=0$ then imply
$W(P_t,Q_t)=0$.

For $m=1$, the same construction starts with $H=A+O(2)$ and changes the
free coefficient $q_{20}$ in \eqref{eq:m1} with nonzero derivative.
Thus the exceptional coefficient is transverse in every integer order.

\section{Proof of the main theorem}

\begin{proof}[Proof of Theorem~\ref{thm:main}]
For $m\ge2$, Proposition~\ref{prop:classification} shows that the
discriminant has type $A_m$ exactly when $\kappa\ne0$.  This is a
finite-jet nonvanishing condition and is therefore open.  If $\kappa=0$,
the analytic family \eqref{eq:deform} consists of genuine fixed-$m$
$W$-congruence germs and, by \eqref{eq:kappat}, has $\kappa_t\ne0$ for
every sufficiently small $t\ne0$.  Thus the $A_m$ locus is dense, and
$d\kappa/dt=-C\ne0$ makes the exceptional locus a smooth hypersurface.

For $m=1$, equation \eqref{eq:m1} and the transverse variation of $q_{20}$
give the same conclusion.  Both signs of the real $A_m$ normal form are
allowed, as required for the unsigned classification.
\end{proof}

\begin{thebibliography}{9}

\bibitem{CraizerGarcia}
M.~Craizer and R.~A.~Garcia,
\emph{Singularities of Weingarten line congruences},
\href{https://arxiv.org/abs/2504.07289}{arXiv:2504.07289}.

\bibitem{Arnold}
V.~I.~Arnold,
\emph{Geometrical Methods in the Theory of Ordinary Differential Equations},
2nd ed., Springer, New York, 1988.

\bibitem{Hormander}
L.~H\"ormander,
\emph{The Analysis of Linear Partial Differential Operators I},
2nd ed., Springer, Berlin, 1990.

\end{thebibliography}
\end{document}
